% ------------------------------------------------------------------------------------
%                           PROGRAMMAZIONE LINEARE - IN DEPTH
% ------------------------------------------------------------------------------------

\chapter{Programmazione Lineare}

Torniamo ora al paradigma della \textbf{programmazione lineare} per
discuterlo più nel dettaglio.

\section{Forma standard e forma canonica}

Abbiamo già visto una possibile forma in cui può presentarsi un problema
di programmazione lineare nel capitolo introduttivo.

Ora, senza perdita di generalità, possiamo assumere che il problema sia espresso 
in \textbf{forma standard}: 

\begin{definition}[FORMA STANDARD]
    \[
        \begin{cases}
            \text{min } c^\top x \\ 
            Ax = b \\ 
            x_i \ge 0
        \end{cases}
    \] 
\end{definition}

Oppure in \textbf{forma canonica}:

\begin{definition}[FORMA CANONICA]
    \[
        \begin{cases}
            \text{min } c^\top x \\ 
            Ax \ge  b \\ 
            x_i \ge 0
        \end{cases}
    \] 
\end{definition}

\begin{remark}
    Il nome di queste forme deriva dal fatto che \textbf{qualsiasi altro problema 
    di programmazione lineare} è \textbf{riconducibile a queste}.
\end{remark}

\subsection{Formulazioni equivalenti}

Possiamo \textbf{passare} da una \textbf{formulazione generica} di un problema
di programmazione lineare ad una \textbf{forma standard (o canonica) equivalente} applicando
le regole di trasformazione seguenti:

\begin{enumerate}
    \item Conversione da \textbf{massimizzazione a minimizzazione}
        \[
            \text{max}(c^\top x) = -\text{min}(-c^\top x)
        \] 
    \item Conversione da \textbf{vincoli} "$\ge $" a vincoli "$=$"
        \[
            a_i^\top x \ge b_i \Longleftrightarrow \begin{cases}
                a_i^\top x - s_i = b_i \implies a_i^\top x = b_i'  \\ 
                s_i \ge 0
            \end{cases}
        \] 
    \item Conversione da \textbf{vincoli} "$\le $" a vincoli "$=$"
        \[
            a_i^\top x \le b_i \Longleftrightarrow \begin{cases}
                a_i^\top x + s_i = b_i \implies a_i^\top x = b_i'  \\ 
                s_i \ge 0
            \end{cases}
        \]
    \item Conversione da \textbf{vincoli} "$=$" a vincoli "$\ge$"
        \[
            a_i^\top x = b_i \Longleftrightarrow \begin{cases}
                a_i^\top x \ge  b_i  \\ 
                -a_i^\top x \ge -b_i
            \end{cases}
        \] 
    \item Conversione da \textbf{variabili libere} a \textbf{non negative}
        \[
            x_i \text{ libera } \Longleftrightarrow \begin{cases}
                x_i = x_i^+ - x_i^- \\ 
                x_i^+,x_i^- \ge 0
            \end{cases}
        \] 
    \item Conversione da \textbf{lower bound generico} a \textbf{variabile non negativa}
        \[
            x_i \ge l_i \Longleftrightarrow \begin{cases}
                x_i = x_i' + l_i \\
                x_i' \ge0
            \end{cases}
        \] 
    \item Conversione da \textbf{upper bound generico} a \textbf{variabile non negativa}
        \[
            x_i \le u_i \Longleftrightarrow \begin{cases}
                x_i = u_i - x_i' \\
                x_i' \ge0
            \end{cases}
        \]
    \item Conversione da \textbf{variabile bounded} a \textbf{variabili non negative}
        \[
            l_i \le  x_i \le u_i \Longleftrightarrow \begin{cases}
                x_i = x_i' + l_i \\
                x_i' + s_i = u_i - l_i \\
                x_i', s_i \ge0
            \end{cases}
        \] 
\end{enumerate}

\begin{warn}
    Il \textbf{passaggio} da una forma all’altra può richiedere una 
    \textbf{variazione del numero di variabili e/o vincoli} del problema; tuttavia
    le trasformazioni inducono un \textbf{blow-up al più polinomiale} della 
    taglia del problema.
\end{warn}

\section{Interpretazione geometrica}

Proviamo ora a dare un'interpretazione geometrica ad un problema di PL.
Per cominciare, diamo le seguenti due definizioni:

\begin{definition}[SEMISPAZIO AFFINE e IPERPIANO]
    Gli insiemi
    \[
        \begin{aligned}
            \{x\in \mathbb{R}^n : a^\top x \le a_0 \} \\ 
            \{x\in \mathbb{R}^n : a^\top x = a_0\}  
        \end{aligned}
    \] 
    Si dicono (rispettivamente) \textbf{semispazio affine} e \textbf{iperpiano}
    \textit{indotti} da $(a,a_0)$.
\end{definition}

\begin{figure}[!ht]
    \centering
    \begin{tikzpicture}[scale=2, x={(-0.6cm,-0.4cm)}, y={(1cm,0cm)}, z={(0cm,1cm)}]
        % Assi coordinati (usiamo un grigio-blu scuro della palette per coerenza)
        \draw[->, nord3] (0,0,0) -- (1.5,0,0) node[anchor=north east]{$x_1$};
        \draw[->, nord3] (0,0,0) -- (0,1.5,0) node[anchor=north west]{$x_2$};
        \draw[->, nord3] (0,0,0) -- (0,0,1.5) node[anchor=south]{$x_3$};

        % Punti per definire il piano (iperpiano)
        \coordinate (A) at (1.5, 0, 0.5);
        \coordinate (B) at (0.5, 1.5, 0.5);
        \coordinate (C) at (-0.5, 1.5, 1.5);
        \coordinate (D) at (0.5, 0, 1.5);

        % Disegno l'iperpiano
        \filldraw[fill=nord9!40, draw=nord10, thick, opacity=0.8] (A) -- (B) -- (C) -- (D) -- cycle;

        % Punto generico x sull'iperpiano
        \coordinate (P) at (0.5, 0.75, 1.0);
        \filldraw[nord12] (P) circle (1pt) node[anchor=north west] {};

        % Vettore normale a
        \draw[->, thick, nord12] (P) -- ++(0.75, 0.3, 0.75) node[anchor=south west] {$a$};

        % Etichetta
        \node[nord10, right] at (B) {$a^\top x = a_0$};
    \end{tikzpicture}
    \caption{Esempio di iperpiano in $\mathbb{R}^3$}
\end{figure}

\begin{definition}[POLIEDRO]
    Si dice \textbf{poliedro} l'\textbf{intersezione} di un
    \textbf{numero finito} di \textbf{semispazi affini} ed \textbf{iperpiani}.
\end{definition}

Notiamo che equazioni di questo tipo sono \textit{proprio quelle che compaiono in un problema di PL}.
L’\textbf{insieme delle soluzioni ammissibili} di un problema di programmazione lineare è quindi, per
definizione, un \textbf{poliedro}; più precisamente, un \textbf{politopo}:

\begin{definition}[POLITOPO]
    Si dice \textbf{politopo} un \textbf{poliedro limitato}
\end{definition}

\begin{figure}[!ht]
    \centering
    \begin{tikzpicture}
        % Definizione dei vertici del poliedro
        \coordinate (V1) at (-0.5, 2);    % Top
        \coordinate (V2) at (1.5, 0.8);   % Top-right
        \coordinate (V3) at (1.5, -1);    % Bottom-right
        \coordinate (V4) at (0, -2);      % Bottom
        \coordinate (V5) at (-1.5, -1.2); % Bottom-left
        \coordinate (V6) at (-1.8, 0.5);  % Top-left

        % Disegno l'area interna del poliedro (con un grigio chiarissimo della palette)
        \fill[nord6] (V1) -- (V2) -- (V3) -- (V4) -- (V5) -- (V6) -- cycle;

        % Etichetta centrale per il poliedro
        \node[nord0] at (0, -0.2) {$P$};

        % Disegno i bordi estesi (iperpiani/rette che definiscono i semispazi)
        \begin{scope}[thick, nord1, shorten <=-12pt, shorten >=-12pt]
            \draw (V1) -- (V2);
            \draw (V2) -- (V3);
            \draw (V3) -- (V4);
            \draw (V4) -- (V5);
            \draw (V5) -- (V6);
            \draw (V6) -- (V1);
        \end{scope}

        % Disegno dei vertici (punti neri)
        \filldraw[nord0] (V1) circle (1.5pt);
        \filldraw[nord0] (V2) circle (1.5pt);
        \filldraw[nord0] (V3) circle (1.5pt);
        \filldraw[nord0] (V4) circle (1.5pt);
        \filldraw[nord0] (V5) circle (1.5pt);
        \filldraw[nord0] (V6) circle (1.5pt);

    \end{tikzpicture}
    \caption{Esempio di politopo in $\mathbb{R}^2$}
    \label{fig:03_02_politopo}
\end{figure}

\subsection{Descrizione interna ed esterna}

Quella che abbiamo visto sopra è la \textbf{descrizione esterna} di un politopo, ovvero
una sua descrizione tramite \textbf{intersezione di vincoli}.

Dalla figura~\ref{fig:03_02_politopo} possiamo tuttavia intuire che è possibile
\textbf{descrivere un politopo} tramite i suoi \textbf{vertici}: si parla in questo
caso di \textbf{descrizione interna}.

\begin{definition}[VERTICE]
    Si dice \textbf{vertice di un poliedro} $P$ un \textbf{punto} $x$ che
    \textbf{non} può essere \textbf{espresso} come \textbf{combinazione convessa stretta} 
    di \textbf{altri due punti del poliedro} stesso. Ovvero, dato il punto
    $x^* \in P$, esso è \textbf{vertice} se:
    \[
        \begin{aligned}
            &\forall x_{1},x_{2} \in P \\ 
            &\nexists \lambda_{1},\lambda_{2}\in(0,1) : \lambda_{1}+\lambda_{2} = 1 \\ 
            &: \lambda_{1}x_{2} + \lambda_{2}x_{2} = x
        \end{aligned}
    \] 
\end{definition}

Ogni politopo ha un \textit{numero finito di vertici} e soddisfa il seguente teorema:

\begin{theorem}
    Ogni \textbf{punto di un politopo} può essere espresso come \textbf{combinazione convessa}
    dei suoi \textbf{vertici} $\{x_{1},\ldots ,x_t\}  $:
    \[
        x\in P \Longleftrightarrow \begin{cases}
            x = \sum_{j=1}^t \lambda_j x_j \\ 
            \sum_{j=1}^t \lambda_j = 1 \\ 
            \lambda_j \ge 0 \text{ , }\forall j=1,\ldots,t
        \end{cases}
    \] 
\end{theorem}

Chiaramente, i vertici sono \textbf{tutti necessari} per descrivere un poliedro.

Ora, per facilitare la compresione della definizione del teorema, riportiamo alcune
note sulle combinazioni convesse:

\begin{remark}[COMBINAZIONE CONVESSA]
    Una combinazione convessa di $t$ punti è una loro \textbf{combinazione lineare} in cui:
    \begin{itemize}
        \item I \textbf{pesi} $\lambda_j$ sono \textbf{non negativi}.
        \item I \textbf{pesi sommano a uno}.
    \end{itemize}

    Con \textbf{due punti}, le loro \textbf{combinazioni convesse} formano il \textbf{segmento 
    di retta} che li unisce. In generale, la combinazione convessa di un insieme di
    punti (non allineati) \textbf{genera la figura convessa} (piena) che li \textbf{"avvolge"},
    chiamata \textbf{involucro convesso} (\textit{convex hull}).
\end{remark}

\begin{exa}[SEGMENTO DI RETTA]
    Poniamo di avere:
    \[
        \begin{aligned}
            &\alpha_{1}x_{1} + \alpha_{2}x_{2} = 0 \\ 
            &\alpha_{1} + \alpha_{2} = 1 \implies \alpha_{2} = 1-\alpha_{1} \\ 
        \end{aligned}
    \] 
    Quindi:
    \[
        \begin{aligned}
            \alpha_{1}x_{1} + (1-\alpha_{1})x_{2} &= \alpha_{1}x_{1} + x_{2} -\alpha_{1}x_{2} \\ 
                                                  &= x_{2} + \alpha_{1}(x_{1}-x_{2})
        \end{aligned}
    \] 
\end{exa}

\subsection{Teorema fondamentale della PL}

L'\textbf{utilità} della \textbf{descrizione interna} di un politopo sta nel
seguente teorema:

\begin{theorem}[TEOREMA FONDAMENTALE DELLA PL]
    Se un problema di programmazione lineare $\text{min}\{cx : x\in P\}$ \textbf{ammette
    ottimo finito}, allora \textbf{esiste un vertice} di $P$ \textbf{ottimo}.
\end{theorem}

Facciamo attenzione a un dettaglio:

\begin{warn}
    I vertici \textbf{non sono gli unci ottimi}, il teorema afferma semplicemente
    che, se l'ottimo della funzione obiettivo è finito, allora \textbf{sicuramente
    un vertice di P sarà ottimo}.
\end{warn}

Per capire meglio, proviamo a considerare gli \textbf{insiemi di livello della funzione
obiettivo}:
\[
    \begin{aligned}
        f(x) &= l \in \mathbb{R} \\ 
        \implies c^\top x &= l
    \end{aligned}
\]
Gli insiemi di livello per la funzione obiettivo di un problema di PL sono degli \textbf{iperpiani}.
Al variare di $l$ otteniamo \textbf{infiniti iperpiani}, \textbf{tutti paralleli}
tra loro (tutti perpendicolari al vettore normale $c$, che rimane invariato).
Notiamo inoltre che $c$ è proprio il \textbf{gradiente} di $f(x)$, per cui:

\begin{remark}[RISOLVERE UN PROBLEMA PL]
    Intuiamo che \textbf{minimizzare} $f(x)$ significa \textbf{spostarsi} negli insiemi di 
    livello \textbf{in direzione} $-c$ (\textit{opposta} al gradiente) fino a trovare
    la \textbf{migliore soluzione} all'\textbf{interno di P}.
\end{remark}

\begin{figure}[!ht]
    \centering
    \begin{tikzpicture}[scale=0.8]
        % Definizione dei vertici del poliedro (soluzioni ammissibili di base)
        \coordinate (V1) at (-0.5, 2);    % Top (soluzione ottima)
        \coordinate (V2) at (1.5, 0.8);   % Top-right
        \coordinate (V3) at (1.5, -1);    % Bottom-right
        \coordinate (V4) at (0, -2);      % Bottom
        \coordinate (V5) at (-1.5, -1.2); % Bottom-left
        \coordinate (V6) at (-1.8, 0.5);  % Top-left

        % Disegno l'area interna del poliedro (regione ammissibile)
        \fill[nord6] (V1) -- (V2) -- (V3) -- (V4) -- (V5) -- (V6) -- cycle;

        % Etichetta centrale per il poliedro
        \node[nord0] at (0, -0.2) {$P$};

        % Disegno i bordi estesi (iperpiani che definiscono i semispazi)
        \begin{scope}[thick, nord1, shorten <=-12pt, shorten >=-12pt]
            \draw (V1) -- (V2);
            \draw (V2) -- (V3);
            \draw (V3) -- (V4);
            \draw (V4) -- (V5);
            \draw (V5) -- (V6);
            \draw (V6) -- (V1);
        \end{scope}

        % Disegno l'iperpiano della funzione obiettivo (curva di livello ottima)
        \draw[thick, nord11] ($(V1) + (-1.5, 0.3)$) -- ($(V1) + (2, -0.4)$) node[right, nord0] {$cx$};

        % Vettore normale c (gradiente della funzione obiettivo)
        \draw[->, thick, nord11] ($(V1) + (1.2, -0.24)$) -- ++(0.24, 1.2) node[above, nord0] {$-c$};

        % Disegno dei vertici non ottimi
        \filldraw[nord0] (V2) circle (1.5pt);
        \filldraw[nord0] (V3) circle (1.5pt);
        \filldraw[nord0] (V4) circle (1.5pt);
        \filldraw[nord0] (V5) circle (1.5pt);
        \filldraw[nord0] (V6) circle (1.5pt);

        % Evidenzio il vertice ottimo sul quale si ferma l'iperpiano
        \filldraw[nord11] (V1) circle (1.5pt);

    \end{tikzpicture}
    \caption{Interpretazione geometrica del teorema fondamentale della PL}
    \label{fig:03_03_InterpretazioneGeomThFondPL}
\end{figure}

La figura~\ref{fig:03_03_InterpretazioneGeomThFondPL} illustra l'interpretazione discussa sopra. 
Ci permette inoltre di intuire perchè \textit{non necessariamente tutti i punti ottimi sono vertici};
se infatti $c$ dovessere essere ortogonale ad uno dei lati del politopo, allora 
\textit{esisterebbe un insieme di livello coincidente con il lato in questione}, dando luogo a
\textit{infiniti punti ottimi} (di cui due sarebbero \textit{comunque vertici}).

Inoltre, quello discusso è proprio l'approccio seguito da uno dei principali algoritmi di
risoluzione per problemi PL:

\begin{definition}[ALGORITMO DEL SIMPLESSO]
    Algoritmo risolutivo per problemi di PL che, \textbf{partendo da un vertice qualsiasi},
    si \textbf{muove in modo greedy} ad un \textbf{vertice vicino non peggiore} ad ogni
    iterazione, \textbf{fino a trovare un vertice ottimo}.
\end{definition}

\section{Esempi di modellizzazione con PL}

Per quanto problemi diversi saranno modellati diversamente, possiamo approcciare 
la fase di modellazione seguendo la seguente \textbf{scaletta generale}:

\begin{remark}[SCALETTA GENERALE]
    La scaletta è la seguente:
    \begin{enumerate}
        \item \textbf{Riconoscere} gli \textbf{insiemi} in gioco.
        \item \textbf{Riconoscere} i \textbf{dati}(\textit{parametri}, detti anche \textbf{attributi}) 
            del problema (\textit{dati in input}).
        \item \textbf{Scegliere} le \textbf{variabili} (\textit{decisioni} del
            problema).
        \item Determinare la \textbf{funzione obiettivo}.
        \item Determinare i \textbf{vincoli} del problema.
    \end{enumerate}
\end{remark}

Di seguito vediamo alcuni esempi in cui il problema è modellato seguendo
questa scaletta.

\subsection{Problema del trasporto}

Il problema è così posto:

\begin{definition}[PROBLEMA DEL TRASPORTO]
    Abbiamo un \textbf{insieme di sorgenti} $\{s_{1},s_{2},\ldots ,s_n\}  $ ed un 
    \textbf{insieme di destinazioni} $\{d_{1},d_{2},\ldots ,d_k\}  $.
    \begin{itemize}
        \item Per ogni coppia \textbf{sorgente-destinazione} $(s_i, m_j)$ 
            abbiamo un \textbf{certo costo} di spedizione $c_{ij}$.
        \item Ogni sorgente $s_i$ ha una \textit{certa quantità di bene}
            (\textbf{disponibilità}) $d_i$.
        \item Ogni destinazione $m_j$ ha una \textit{certa domanda} 
            di prodotto (\textit{richiesta}) $r_j$.
    \end{itemize}

    Vogliamo ideare il \textbf{piano di spedizione} che \textbf{minimizza
    il costo totale di spedizione}.
\end{definition}

\begin{figure}[!ht]
    \centering
    \begin{tikzpicture}[scale=1.5,
        % Stili dei nodi
        source/.style={circle, draw=nord3, fill=nord4, thick, minimum size=8mm, inner sep=0pt},
        dest/.style={circle, draw=nord3, fill=nord6, thick, minimum size=8mm, inner sep=0pt},
        % Stile degli archi
        edge/.style={->, thick, nord3, shorten >=2pt, shorten <=2pt}
        ]

        % --- Posizionamento dei Nodi ---
        % Nodi Sorgente (A sinistra)
        \node[source] (A) at (0, 1) {A};
        \node[source] (B) at (0, -1) {B};

        % Etichette disponibilità (d_i)
        \node[] at (-0.9,1) {$d_A = 250$};
        \node[] at (-0.9,-1) {$d_B = 350$};

        % Nodi Destinazione (A destra)
        \node[dest] (1) at (3, 1.5) {1};
        \node[dest] (2) at (3, 0) {2};
        \node[dest] (3) at (3, -1.5) {3};

        % Etichette richieste (r_j)
        \node[] at (3.9, 1.5) {$r_1 = 170$};
        \node[] at (3.9, 0) {$r_2 = 100$};
        \node[] at (3.9, -1.5) {$r_3 = 80$};

        % --- Archi (Costi c_{ij}) ---
        % Da A a tutte le destinazioni
        \draw[edge] (A) -- (1);
        \draw[edge] (A) -- (2);
        \draw[edge] (A) -- (3);

        % Da B a tutte le destinazioni
        \draw[edge] (B) -- (1);
        \draw[edge] (B) -- (2);
        \draw[edge] (B) -- (3);

        % --- Etichette Gruppi ---
        \node[nord0, font=\bfseries] at (0, 2.2) {Sorgenti (I)};
        \node[nord0, font=\bfseries] at (3, 2.2) {Destinazioni (J)};

    \end{tikzpicture}
    \caption{Esempio di problema del trasporto}
\end{figure}

\begin{remark}
    La rappresentazione del problema è data da un \textbf{grafo bipartito}.
\end{remark}

Nell'esempio in figura abbiamo a che fare con i seguenti insiemi:
\[
    \begin{aligned}
        I &: \{\text{sorgenti}\} = \{A,B\} \\ 
        J &: \{\text{destinazioni}\}  = \{1,2,3\}  
    \end{aligned}
\]
Abbiamo inoltre i seguenti \textit{dati di input} (\textbf{parametri}): 
\begin{itemize}
    \item \textbf{Disponibilità}: $d_A=250$, $d_B=350$
    \item \textbf{Richieste}: $r_{1}=170$, $r_{2}=100$, $r_{3}=80$.
    \item \textbf{Costi unitari} di spedizione: $c_{ij}$.
\end{itemize}

E scegliamo di utilizzare come \textbf{variabili}:
\[
    x_{ij}\in \mathbb{R}^+ = \text{ q.tà di prodotto da spedire da }i \text{ a }j \text{ , }
    (\forall i\in I \text{ , }\forall j \in J)
\]

Dove con $\mathbb{R}^+$ intendiamo $x\in \mathbb{R}, x\ge 0$ (quindi 0 incluso).

\begin{remark}
    La \textit{descrizione semantica} delle variabili scelte \textbf{non}
    \textit{fa parte del modello}: si tratta semplicemente di una \textit{descrizione
    utile alla compresione del modello}, ma, ovviamente, l'\textit{algoritmo
    risolutivo non ne è a conoscenza}.
\end{remark}

Dato che vogliamo \textbf{minimizzare il costo totale di spedizione}, ha senso 
scegliere la seguente \textbf{funzione obiettivo}:
\[
    f(x) = \sum_{i\in I} \sum_{j\in J} c_{ij}x_{ij}
\]

Infine, per quanto riguarda i vincoli, possiamo considerare i seguenti:
\begin{enumerate}
    \item Le \textbf{richieste} delle destinazioni devono essere \textbf{soddisfatte}: 
        $\sum_{i\in I} x_{ij} \ge r_j \text{ } \forall j\in J$
    \item \textbf{Non} bisogna \textbf{eccedere le disponibilità}: 
        $\sum_{j\in J}x_{ij} \le d_i \text{ } \forall i\in I$
\end{enumerate}

Terminato questo processo di modellazione, il nostro modello per il problema 
in esempio risulta quindi essere:
\[
    \begin{cases}
        \text{min } \sum_{i\in I} \sum_{j\in J} c_{ij}x_{ij} \\ \\  
        \sum_{i\in I} x_{ij} \ge r_j & \forall j\in J \\ \\ 
        \sum_{j\in J}x_{ij} \le d_i & \forall i\in I \\ \\
        x_{ij} \in \mathbb{R}^+ & \forall i \in I \text{ , }\forall j\in J
    \end{cases}
\]

\subsection{Problema della dieta}

Il problema è così posto:

\begin{definition}[PORBLEMA DELLA DIETA]
    Abbiamo:
    \begin{itemize}
        \item Un \textbf{insieme di alimenti} $J=\{\text{alimenti}\}  $ di una \textit{possibile dieta} 
        (per animali da fattoria da sfamare), con \textbf{costo unitario} $c_j$. 
        \item Una \textbf{lista di sostanze nutritive} $I=\{\text{sostanze nutritive}\}$, 
        ciasuna con \textbf{quantità minima} $b_i$. 
        \item Dei valori $a_{ij}$ che rappresentano la \textbf{quantità} di \textbf{nutriente} 
        $i$ che si trova \textbf{nell'alimento} $j$.
    \end{itemize}
\end{definition}

Scegliamo come \textbf{variabili}:
\[
    x_j \in \mathbb{R}^+ = \text{q.tà di elemento j nella dieta }\forall j\in J
\]

Intuitivamente, il modello sarà:
\[
    \begin{cases}
        \text{min } \sum_{j\in J} c_j x_j \\ \\ 
        \sum_{j\in J} a_{ij}x_j \ge b_i & \forall i\in I \\ \\ 
        x_j \ge 0 & \forall j\in J
    \end{cases}
\]

\begin{remark}
    Vogliamo \textbf{minimizzare} il \textbf{costo della dieta} di un animale, rispettando
    comunque i \textbf{fabbisogni minimi} $b_i$ per \textbf{ogni nutriente}.
\end{remark}

\section{Dualità}

Ricordiamo che \textbf{risolvere un problema} di ottimizzazione \textbf{non} significa 
\textbf{solamente trovare una soluzione} ottima, è anche necessario \textbf{dimostrare}
che la soluzione \textbf{sia effettivamente ottima}.

Tenendo questo a mente, riprendiamo il concetto di \textbf{dualità}. 
Abbiamo visto che un problema di ottimizzazione può essere visto come:
\begin{enumerate}
    \item Un problema (\textbf{primale}) di \textbf{ricerca}: cerchiamo 
        \textit{soluzioni via via migliori} generando \textbf{nuovi upper bound}.
    \item Un problema (\textbf{duale}) di \textbf{inferenza}: usiamo i \textit{vincoli}
        per generare \textbf{lower bound} via via migliori.
\end{enumerate}

Eseguendo \textbf{in parallelo} la risoluzione del \textbf{primario} e del \textbf{duale}
possiamo \textbf{trovare una soluzione} e anche \textbf{dimostrarne l'ottimalità}.

\begin{figure}[!ht]
    \centering
    \begin{tikzpicture}[scale=1.2, >=stealth]
        % Assi coordinati (in grigio scuro per risaltare)
        \draw[->, ultra thick, nord0] (-0.5, 0) -- (6.5, 0);
        \draw[->, ultra thick, nord0] (0, -0.6) -- (0, 3.5);

        % Livello della soluzione ottima (sol)
        \draw[dashed, thick, nord0] (0, 1.5) -- (6, 1.5);
        \node[left, nord0, font=\bfseries\Large] at (-0.2, 1.5) {sol};

        % Curva Primale P (Upper bound che scende, colore verde Nord)
        \draw[ultra thick, nord14] (0, 3.1) to[out=-40, in=175] (6, 1.55);
        \node[above, nord14, font=\bfseries\Large] at (5.5, 1.8) {P};

        % Curva Duale D (Lower bound che sale, colore rosso Nord)
        \draw[ultra thick, nord11] (0, 0) to[out=30, in=185] (6, 1.45);
        \node[below, nord11, font=\bfseries\Large] at (5.5, 1.2) {D};

    \end{tikzpicture}
    \caption{Risoluzione parallela del problema primario e del duale}
\end{figure}

\subsection{Formulazione matematica del duale}

Vediamo ora più nel dettaglio in cosa consiste il problema \textbf{duale}.

Dato un problema di ottimizzazione $P$ del tipo
\[
    P: \begin{cases}
        \text{min } c^\top x \\ 
        Ax \ge b
    \end{cases}
\]
possiamo intuire dalla figura e da quanto detto sopra che il suo \textbf{duale} 
consiste nel \textbf{determinare un lower bound} $c_{0}$ per $c^\top x$ e 
può essere formulato come segue:

\begin{remark}[FORMULAZIONE DEL DUALE]
    Risolvere il problema \textbf{duale} significa \textbf{cercare disuguaglianze}
    del tipo:
    \[
        c^\top x \ge c_{0}
    \] 
    In modo che $c_{0}$ sia \textbf{massimo} (lower bound più stringente possibile).
    Cioè:
    \[
        \begin{cases}
            \text{max } c_{0} \\ 
            c^\top x \ge c_{0}
        \end{cases}
    \]
\end{remark}

A questo punto dobbiamo però precisare che la disuguaglianza che stiamo cercando
deve \textbf{essere valida}. Formalizziamo allora questo concetto:

\begin{definition}[DISUGUAGLIANZA VALIDA]
    Data una \textbf{regione ammissibile} $P$ ed una \textbf{disuguaglianza}
    $\alpha^\top x \ge \beta$, diciamo che questa è \textbf{valida} in $P$
    \textbf{se e solo se} vale:
    \[
        \alpha^\top x \ge \beta \text{ } \forall x \in P
    \] 
\end{definition}

Verificare che la \textbf{disuguaglianza sia effettivamente rispettata} $\forall x\in P$ è 
chiaramente \textbf{infattibile} dato che esistono \textbf{infiniti punti} in $P$; per 
fortuna, però, possiamo sfruttare il seguente teorema:

\begin{theorem}
    Se una \textbf{disuguaglianza} è \textbf{valida per tutti i vertici} di un politopo
    $P$, allora è valida \textbf{anche per tutti} gli 
    \textbf{altri punti} di $P$.
\end{theorem}

Del resto, se $P$ è identificato dai vertici $x_{1},x_{2},\ldots ,x_t$, sappiamo
che un punto $x\in P$ può essere scritto come:
\[
    x = \lambda_{1}x_{1} + \ldots \lambda_t x_t
\]

Con $\sum \lambda_j = 1$ e $\lambda_j\ge 0$, $j=1,\ldots ,t$. Abbiamo quindi:
\[
    \begin{aligned}
        \alpha^\top x &= \alpha^\top (\lambda_{1}x_{1} + \ldots \lambda_t x_t) \\ 
              &= \lambda_{1}(\alpha^\top x_{1}) + \ldots + \lambda_t(\alpha^\top x_t)
    \end{aligned}
\]

Supponiamo ora che la diseguaglianza sia \textbf{verificata per ogni vertice}:
\[
    \begin{aligned}
        \alpha^\top x_{1} &\ge \beta \\ 
        \alpha^\top x_{2} &\ge \beta \\ 
               &\ldots \\ 
        \alpha^\top x_t &\ge \beta
    \end{aligned}
\]

Allora risulta:
\[
    \begin{aligned}
        \alpha^\top x &\ge \lambda_{1}(\beta) + \ldots + \lambda_t(\beta) \\
                      &\Downarrow \\ 
        \alpha^\top x &\ge \beta(\lambda_{1}+\ldots +\lambda_t) = \beta 
    \end{aligned}
\]

Esattamente come volevamo $\square$.

\subsection{Approccio e lemma di Farkas}

Abbiamo quindi capito che risolvere il problema duale significa 
cercare il \textbf{massimo} $c_{0}$ tale che $c^\top x \ge c_{0}$ sia \textbf{valida} per
il problema $P$ che stiamo risolvendo. Abbiamo inoltre capito \textit{come}
verificare se una disuguaglianza è \textit{valida o meno}, ma non abbiamo
ancora discusso \textbf{come trovare} $c_{0}$.

Nella pratica, la ricerca di $c_{0}$ si traduce nello \textbf{sfruttare i vincoli} $Ax\ge b$ 
per \textit{dedurre} (\textbf{inferenza}) dei 
\textbf{lower bound stringenti} $c^\top x \ge c_{0}$.

Intuitivamente, possiamo pensare di ottenere tali lower bound 
\textbf{moltiplicando la disequazione} data dai \textbf{vincoli} 
per un \textbf{vettore} $u\ge 0$ (componenti \textit{non negative}):
\[
    \underbrace{ u^\top A }_{c} x \ge \underbrace{ u^\top b }_{c_{0}}
\]

Il problema diventa quindi:
\[
    \begin{cases}
        \text{max } c_{0} \\ 
        c^\top x \ge c_{0} \\
        u \ge 0 
    \end{cases}
\]

\begin{remark}[CONDIZIONI SU $u$]
    Siccome il nostro problema \textbf{originale} aveva come 
    \textbf{vincoli} $Ax\ge b$, dobbiamo \textbf{imporre} $u\ge 0$ per
    \textbf{non modificare il verso} della disuguaglianza.

    Se \textbf{invece} i vincoli fossero stati della forma
    $Ax=b$, allora $u$ sarebbe stato \textbf{libero}.
\end{remark}

Ora, chiaramente \textbf{non possiamo procedere alla cieca}: $u$ dovrà
essere \textbf{scelto con criterio} affinchè la disuguaglianza
\textbf{risulti valida}. Per questo, ci viene in aiuto il seguente lemma:

\begin{theorem}[LEMMA DI FARKAS]
    Una \textbf{disuguaglianza} $c^\top x \ge c_{0}$ è \textbf{valida}
    per $P$ \textbf{se e solo se}:
    \[
        \exists u \in \mathbb{R}_+^m \text{ t.c. } 
        c^\top \ge u^\top A \text{ e } c_{0} \le u^\top b
    \] 
\end{theorem}

Il lemma di Farkas ci permette di riscrivere il problema come segue:
\[
    \begin{cases}
        \text{max } c_{0} \\ 
        c^\top \ge u^\top A \\
        c_{0} \le u^\top b \\ 
        u \ge 0
    \end{cases}
\]

\begin{remark}[CAMBIO DI PROSPETTIVA] 
    Notiamo che:
    \begin{itemize}
        \item Noi \textbf{vogliamo massimizzare} $c_{0}$.
        \item Contemporaneamente, il \textbf{vincolo} $c_{0} \le u^\top b$ 
            impone un \textit{upper bound} a $c_{0}$. 
    \end{itemize}
    Chiaramente, la \textbf{scelta ottimale} è \textbf{far coincidere}
    $c_{0}$ con il suo \textit{upper bound}, per cui imponiamo
    \[
        c_{0} \coloneqq u^\top b
    \] 
    E scegliamo di \textbf{massimizzare questa quantità}.
\end{remark}

Alla luce di ciò, possiamo riscrivere il problema nella seguente forma:
\[
    \begin{cases}
        \text{max } u^\top b \\ 
        u^\top A \le c^\top \\ 
        u \ge 0
    \end{cases}
\]

\begin{remark}
    Osserviamo che:
    \begin{enumerate}
        \item Il \textbf{duale} di un problema di PL è un \textbf{altro problema di PL}.
        \item I \textbf{ruoli} dei \textbf{vettori} $c$ e $b$ si sono \textbf{scambiati}.
            Ciò significa che:
            \begin{enumerate}
                \item Il numero di \textbf{vincoli del duale} è uguale
                    al numero di \textbf{variabili del primale}.
                \item Il numero di \textbf{variabili del duale} è uguale
                    al numeri di \textbf{vincoli del primale}.
            \end{enumerate}
    \end{enumerate}
\end{remark}

Il secondo punto è schematizzato in figura.

\begin{figure}[!ht]
    \centering
    \begin{tikzpicture}[
        scale=1,
        % Stili personalizzati per i blocchi
        matrix block/.style={draw=nord0, fill=nord5, rounded corners=2pt, 
        ultra thick, align=center, font=\Huge\bfseries},
        vector block/.style={draw=nord0, fill=nord4, rounded corners=2pt, 
        ultra thick, align=center, font=\Large\bfseries},
        label/.style={nord0, font=\Huge\bfseries}
        ]

        % ==========================================
        % PROBLEMA PRIMALE (P) - A sinistra
        % ==========================================

        % Matrice A (larga)
        \draw[matrix block] (-2, -1.2) rectangle (2, 1.2) node[midway] {$A$};

        % Vettore c (orizzontale in alto)
        \draw[vector block] (-2, 1.35) rectangle (2, 1.95) node[midway] {$c$};

        % Vettore b (verticale a destra)
        \draw[vector block] (2.15, -1.2) rectangle (2.75, 1.2) node[midway] {$b$};

        % Etichetta P in basso
        \node[label] at (0, -2) {$P$};


        % ==========================================
        % FRECCIA CENTRALE
        % ==========================================
        \draw[<->, line width=2.5pt, nord3] (3.5, 0) -- (4.5, 0);


        % ==========================================
        % PROBLEMA DUALE (D) - A destra
        % ==========================================

        % Matrice A^T (trasposta, quindi alta e stretta)
        \draw[matrix block] (6.8, -2) rectangle (9.2, 2) node[midway] {$A^T$};

        % Vettore b (diventa orizzontale in alto)
        \draw[vector block] (6.8, 2.15) rectangle (9.2, 2.75) node[midway] {$b$};

        % Vettore c (diventa verticale a destra)
        \draw[vector block] (9.35, -2) rectangle (9.95, 2) node[midway] {$c$};

        % Etichetta D in basso
        \node[label] at (8, -2.8) {$D$};

    \end{tikzpicture}
    \caption{I ruoli delle variabili e dei vincoli nel duale sono scambiati}
\end{figure}

\subsection{Dualità debole e dualità forte}

Abbiamo visto che il problema \textbf{duale} può essere usato per
produrre dei \textbf{lower bound} per il problema \textbf{primale}.
Ciò è formalizzato dal concetto di \textbf{dualità debole}:

\begin{prop}[DUALITA' DEBOLE]
    \[
        \begin{aligned}
            \forall x\in P &\text{ , } \forall u \in D \\ 
            u^\top b  &\le c^\top x
        \end{aligned}
    \] 
\end{prop}

Cioè, per \textbf{qualsiasi sol.} $x$ del problema \textbf{primale} $P$ e
per \textbf{qualsiasi sol.} $u$ del problema \textbf{duale} $D$, 
il valore della \textbf{funzione obiettivo del duale} ($u^\top b$) sarà 
\textbf{sempre minore o uguale} al valore della \textbf{funzione obiettivo del primale}.

La dimostrazione parte dalle condizioni date dal lemme di Farkas ed è, molto 
banalmente, la seguente:
\[
    c^\top \ge u^\top A \implies c^\top x \ge u^\top A x \ge u^\top b
\]

Inoltre, è importantissimo anche il concetto di \textbf{dualità forte}:

\begin{prop}[DUALITA' FORTE]
    Il \textbf{valore ottimo del duale} ($d^*$) \textbf{coincide} con il
    \textbf{valore ottimo del primale} ($p^*$) \textit{SE E SOLO SE}
    \textbf{entrambi ammettono ottimi finiti}.
\end{prop}

Possiamo interpretare questi due concetti come segue:

\begin{remark}[DUALITA' DEBOLE VS FORTE]
    La dualità forte è un concetto appunto \textit{più forte} perchè \textit{non solo}
    formalizza la \textit{ricerca dei lower bound} per la soluzione, ma ci dice anche 
    \textbf{quando soluzione del primale} e \textbf{soluzione del duale coincidono}.
\end{remark}

Inoltre, questi concetti ci permettono di affermare che \textbf{solo alcune casistiche}
riguardanti i \textbf{due problemi} sono effettivamente \textbf{ammissibili}.
In figura sono schematizzate le situazioni ammissibili (contrassegnate da una spunta), e
quelle \textit{impossibili} (contrassegnate da una croce):

\begin{figure}[!ht]
    \centering
    \begin{tikzpicture}[scale=1.8]
        \newcommand{\cmark}{\textcolor{nord14}{\Huge $\checkmark$}}
        \newcommand{\xmark}{\textcolor{nord11}{\Huge $\times$}}

        % --- Linee della Griglia ---
        % Orizzontali
        \draw[ultra thick, nord3] (0, 3) -- (4, 3);
        \draw[ultra thick, nord3] (0, 2) -- (4, 2);
        \draw[ultra thick, nord3] (0, 1) -- (4, 1);
        \draw[ultra thick, nord3] (0, 0) -- (4, 0);
        \draw[ultra thick, nord3] (0, 4) -- (4, 4);

        % Verticali
        \draw[ultra thick, nord3] (0, 0) -- (0, 4);
        \draw[ultra thick, nord3] (1, 0) -- (1, 4);
        \draw[ultra thick, nord3] (2, 0) -- (2, 4);
        \draw[ultra thick, nord3] (3, 0) -- (3, 4);
        \draw[ultra thick, nord3] (4, 0) -- (4, 4);

        % Diagonale nella cella in alto a sinistra (P/D)
        \draw[ultra thick, nord3] (0, 4) -- (1, 3);
        \node[nord0, font=\Large\bfseries] at (0.3, 3.25) {P};
        \node[nord0, font=\Large\bfseries] at (0.7, 3.75) {D};

        % --- Intestazioni ---
        % Colonne (Duale)
        \node[nord0, font=\Large\bfseries, align=center] at (1.5, 3.5) {Imp.};
        \node[nord0, font=\Large\bfseries, align=center] at (2.5, 3.5) {Sol.\\Finita};
        \node[nord0, font=\Large\bfseries, align=center] at (3.5, 3.5) {Illim.};

        % Righe (Primale)
        \node[nord0, font=\Large\bfseries, align=center] at (0.5, 2.5) {Imp.};
        \node[nord0, font=\Large\bfseries, align=center] at (0.5, 1.5) {Sol.\\Finita};
        \node[nord0, font=\Large\bfseries, align=center] at (0.5, 0.5) {Illim.};

        % --- Contenuto della Tabella ---
        % Riga 1 (P: Imp.)
        \node at (1.5, 2.5) {\cmark};
        \node at (2.5, 2.5) {\xmark};
        \node at (3.5, 2.5) {\cmark};

        % Riga 2 (P: Sol. Finita)
        \node at (1.5, 1.5) {\xmark};
        \node at (2.5, 1.5) {\cmark};
        \node at (3.5, 1.5) {\xmark};

        % Riga 3 (P: Illim.)
        \node at (1.5, 0.5) {\cmark};
        \node at (2.5, 0.5) {\xmark};
        \node at (3.5, 0.5) {\xmark};

        % --- Evidenziazione "Casi Interessanti" ---
        \draw[nord13, ultra thick, dashed, rotate around={45:(2.5, 1.5)}] (2.5, 1.5) ellipse (1.9cm and 0.4cm);
    \end{tikzpicture}
    \caption{Casistiche ammesse dai concetti di dualità.}
\end{figure}

In figura sono evidenziati alcuni \textbf{casi interessanti} (contorno tratteggiato), cioè 
quelli che più spesso si manifestano negli esercizi. 

In ogni caso, elenchiamo i casi possibili e la motivazione che sta dietro:

\begin{remark}
    I seguenti casi sono possibili:
    \begin{itemize}
        \item \textbf{OTTIMO P - OTTIMO D} - Caso descritto dalla \textbf{dualità forte}.
        \item \textbf{INFINITO P - IMPOSSIBILE D} - Conseguenza della \textbf{dualità debole}.
        \item \textbf{IMPOSSIBILE P - INFINITO D} - Conseguenza della \textbf{dualità debole}.
        \item \textbf{IMPOSSIBILE P - IMPOSSIBILE D} - Conseguenza della \textbf{dualità debole}.
    \end{itemize}
\end{remark}

E' facile intuire che queste situazioni sono dirette conseguenze del fatto che 
la \textbf{soluzione del primale} rappresenta un \textbf{upper bound} per il problema, 
mentre quella del \textbf{duale} rappresenta un \textbf{lower bound}.

Inoltre, quando utilizziamo un software per risolvere un problema di PL, è bene
tenere a mente:

\begin{remark}[SOTTO-PRODOTTO DELL'ALGORITMO]
    Quando usiamo un \textbf{algoritmo} (e.g. \textit{algoritmo del simplesso}) per
    risolvere un problema, questo ci fornirà una \textbf{coppia di soluzioni}
    \[
        (x^*, u^*)
    \] 
    Che sono, rispettivamente, \textbf{soluzione del primale} e \textbf{soluzione del duale}.
\end{remark}

Ciò è una \textbf{conseguenza} del fatto che, come detto in precedenza, 
\textbf{non basta trovare una soluzione} per dire di aver \textbf{risolto il problema}, 
ma è anche necessario fornire la \textbf{dimostrazione di ottimalità}. 
Tale dimostrazione è data proprio dalla soluzione del duale.

\subsection{Esempio}

A questo punto, consideriamo alcuni semplici esempi per capire concretamente 
come operare con il problema duale.

\begin{exa}[ESEMPIO DUALITA']
    Consideriamo il seguente esempio:
    \[
        \begin{cases}
            \text{min } 4x_{1} - 2x_{2} + x_{3} + 5x_{4} & \text{f. obiettivo} \\ 
            x_{1} + x_{2} \le 1 & \text{vincolo} \\ 
            4x_{2} - 3x_{3} \ge 2 & \text{vincolo} \\ 
            2x_{1} + x_{3} - x_{4} = 3 & \text{vincolo} \\ 
            x_{1},x_{2} \ge 0 & \text{dom. variabili} \\ 
            x_{3} \le 0 & \text{dom. variabili} \\
            x_{4} \text{ libera} & \text{dom. variabili}
        \end{cases}
    \] 
\end{exa}

Osserviamo subito che il problema ha \textbf{4 variabili} e \textbf{3 vincoli}, per cui
ci aspettiamo che il \textbf{duale} abbia \textbf{3 variabili} e \textbf{4 vincoli}.

In particolare, abbiamo:
\[
    \begin{aligned}
        x = \begin{bmatrix} x_{1} \\ x_{2} \\ x_{3} \\ x_{4} \end{bmatrix} & &
        c^\top = [4,-2,1,5] & & \\ \\
        b = \begin{bmatrix} 1 \\ 2 \\ 3 \end{bmatrix} & &
        A = \begin{bmatrix} 1 & 1 & 0 & 0 \\ 0 & 4 & -3 & 0 \\ 2 & 0 & 1 & -1 \end{bmatrix} 
    \end{aligned}
\]

Per cominciare, determiniamo le condizioni su $u^\top = (u_{1},u_{2},u_{3})$ a 
partire dai vincoli:
\[
    \begin{aligned}
        x_{1} + x_{2} \le 1 &\implies u_{1} \le 0 \\ 
        4x_{2} - 3x_{3} \ge  2 &\implies u_{2} \ge 0 \\
        2x_{1} + x_{3} - x_{4} = 3 &\implies u_{3} \text{ libero}
    \end{aligned}
\]

Inoltre, dai \textbf{termini noti dei vincoli} ricaviamo
la \textbf{funzione obiettivo del duale}:
\[
    u^\top b = u_{1} + 2u_{2} + 3u_{3}
\]
Che, ovviamente, vogliamo \textbf{massimizzare}.

Ora, per ottenere i vincoli dobbiamo considerare $u^\top A$ e $c^\top$. 
Per comodità, consideriamo invece $A^\top u$ e $c$:
\[
    \begin{bmatrix} 
        1 & 0 & 2 \\ 
        1 & 4 & 0 \\ 
        0 &-3 & 1 \\ 
        0 & 0 & -1
    \end{bmatrix} \begin{bmatrix} u_{1} \\ u_{2} \\ u_{3} \end{bmatrix} \text{ , } 
    \begin{bmatrix} 4 \\ -2 \\ 1 \\ 5 \end{bmatrix} 
\]

Abbiamo così \textit{determinato anche i vincoli}, ed il \textbf{problema duale} 
risulta quindi essere:
\[
    \begin{cases}
        \text{max } u_{1} + 2u_{2} +3u_{3} \\ 
        u_{1} + 2u_{3} \le 4 \\ 
        u_{1} + 4u_{2} \le -2 \\ 
        -3u_{2} + u_{3} \ge 1 \\ 
        -u_{3} = 5 \\ 
        u_{1} \le 0 \\ 
        u_{2} \ge 0 \\
        u_{3} \text{ libero}
    \end{cases}
\]

\begin{remark}[VERSI DELLE DISEQUAZIONI NEI VINCOLI DUALI]
    I \textbf{versi delle disequazioni} nei \textbf{vincoli} del \textbf{problema duale}
    sono ottenuti a partire dalle \textbf{condizioni sui domini} delle
    \textbf{variabili del primale}. In particolare:
    \begin{itemize}
        \item Un dominio del tipo $x_i\ge 0$ produce il 
            \textbf{vincolo duale} $\le $.
        \item Un dominio del tipo $x_i \le 0$ produce il
            \textbf{vincolo duale} $\ge $.
        \item Un dominio del tipo $x_i \text{ libero}$ produce il 
            \textbf{vincolo duale} $=$.
        \item Un dominio del tipo $x_i = 0$ produce il
            \textbf{vincolo duale} $\text{libero}$.
    \end{itemize}
\end{remark}
